% ==============================================================================
% РАЗДЕЛ 2.4 (DISS-005-B1292 .. DISS-005-B1349)
% Анализ работы реле направления мощности с применением ТТ и КР
% ==============================================================================

% DISS-005-B1292
\section{Анализ работы реле направления мощности с применением ТТ и КР}
\label{sec:ch02_rnm_tt_kr_comparison}

% DISS-005-B1293
Для проведения исследований функционирования реле направления мощности в граничных случаях в среде Simulink была собрана модель подстанция 230/11 кВ с трансформаторами Т1 и Т2 мощностью 63~МВА и синхронной двигательной нагрузкой (СД) SM-1 и SM-2 мощностью 6~МВт. Помимо СД, на рассматриваемых секциях 3 и 4 так же присутствует стандартная трансформаторная (пассивная) нагрузка (Рисунок~\ref{fig:ch02_simulink_substation_model}).

% DISS-005-B1294
% DISS-005-B1295
\begin{figure}[htbp]
  \centering
  \includegraphics[width=0.88\textwidth]{image68.png}
  \caption{Модель исследуемой подстанции 230/11кВ в среде Simulink.}
  \label{fig:ch02_simulink_substation_model}
\end{figure}

% DISS-005-B1296
В сетях внешнего и внутреннего электроснабжения предприятий может произойти множество различных аварийных событий, в большинстве из которых к классические и адаптивные РНМ функционируют согласно заявленному алгоритму. Поэтому мы рассмотрим случаи, которые возникали в подобных реальных сетях и приводили либо к ложному срабатыванию, либо к задержкам срабатывания РНМ.

% DISS-005-B1297
Используется формула РНМ на основе работ Жукова В.А. [25] с тем отличием, что уставка задаётся не относительно действительной составляющей тока ($\mathrm{Re}(I)$), а строится отдельная характеристики срабатывания.

% DISS-005-B1298
Ниже представлены формулы РНМ, основанного на векторах прямой последовательности (ПП) и векторах фазного тока (2.7):

% DISS-005-B1299
% DISS-005-B1300
% DISS-005-B1301
\begin{equation}
\label{eq:ch02_pdr_tt}
\left\{
\begin{aligned}
\text{ПП: } & \dot{PDR1} = \frac{(\dot{U}_{1,B1} + k \cdot \dot{U}_{1,B2}) \cdot \bar{\dot{I}}_{1,B1}}{|\dot{U}_{1,B1} + k \cdot \dot{U}_{1,B2}|} \\
\text{Фазный: } & \dot{PDR2}_A = \frac{(\dot{U}_{BC,B1} + k \cdot \dot{U}_{BC,B2}) \cdot \bar{\dot{I}}_{A,B1} \cdot e^{-90^\circ}}{|\dot{U}_{BC,B1} + k \cdot \dot{U}_{BC,B2}|}
\end{aligned}
\right.
\end{equation}

% DISS-005-B1302

% DISS-005-B1303
Формулы РНМ для КР (2.8):

% DISS-005-B1304
% DISS-005-B1305
% DISS-005-B1306
\begin{equation}
\label{eq:ch02_pdr_kr}
\left\{
\begin{aligned}
\text{ПП: } & \dot{PDR3} = \frac{(\dot{U}_{1,B1} + k \cdot \dot{U}_{1,B2}) \cdot \bar{\dot{I}}_{1,B1} \cdot e^{-90^\circ}}{|\dot{U}_{1,B1} + k \cdot \dot{U}_{1,B2}|} \\
\text{Фазный: } & \dot{PDR4}_A = \frac{(\dot{U}_{BC,B1} + k \cdot \dot{U}_{BC,B2}) \cdot \bar{\dot{I}}_{A,B1} \cdot e^{-180^\circ}}{|\dot{U}_{BC,B1} + k \cdot \dot{U}_{BC,B2}|}
\end{aligned}
\right.
\end{equation}

% DISS-005-B1307

% DISS-005-B1308
где $\dot{PDR}$~-- итоговый вектор РНМ прямой последовательности или фазный, $\dot{U}_{1,B1}$~-- вектора напряжения прямой последовательности или перпендикулярного линейного напряжения (BC для фазы А) для данной секции шин (B1) или смежной (B2), $k$~-- коэффициент связи, $\bar{\dot{I}}_{1,B1}$~-- сопряжённый вектор тока прямой последовательности или исследуемой фазы.

% DISS-005-B1309
Усреднение сигналов с ТТ и КР производится за счёт сложения выходных комплексных значений векторов РНМ, формулы (2.7) и (2.8) по выражению идентичному формуле (2.4). Гибридный же алгоритм на первом этапе предлагается осуществлять по следующему алгоритму:

% DISS-005-B1310
Рассчитывать производную с датчиков ТТ.

% DISS-005-B1311
Находить разность математической производной тока (с ТТ) и физической (с КР). Для удобства сравнения сигналы были приведены к масштабу токов делением значений на $\omega$.

% DISS-005-B1312
По уровню разницы производных рассчитывать коэффициент гибридной связи ($k$). Для случаев погрешности менее $d = 10\%$ от $I_{\text{ном}}$, предлагается брать сигнал с ТТ ($k = 0$), в промежуточном диапазоне $d = 10 \div 20\%$ осуществлять плавный переход ($k = \frac{d-10}{10}$), а для случая с большей погрешностью использовать сигнал с КР ($k = 1$). 

% DISS-005-B1313
Для каждой фазы производить расчёт гибридного сигнала по формуле (2.9).

% DISS-005-B1314
% DISS-005-B1315
% DISS-005-B1316
\begin{equation}
\label{eq:ch02_hybrid_signal}
\dot{I}_{\text{гибрид.}} = (1 - k) \cdot \dot{I}_{\text{ТТ}} + k \cdot (-j) \cdot \dot{I}_{\text{КР}}
\end{equation}

% DISS-005-B1317
На основе скорректированных сигналов рассчитывать вектор РНМ по классической формуле (2.7).

% DISS-005-B1318
Пороги гибридного алгоритма могут быть уточнены в дальнейших исследованиях. На начальном этапе значение $k$ принимается равным допустимой погрешности классических ТТ, используемых для РЗА (10\%). Отметим, что применением гибридного сигнала позволяет исключить насыщение ТТ выше заданного порога, но никак не сказывается на математических особенностях влияния апериодической составляющей.

% DISS-005-B1319
Первый такой режим~-- это удалённое однофазное (А) короткое замыкание в сети 230~кВ (точка К-1, Рисунок~\ref{fig:ch02_simulink_substation_model}). При возникновении подобного аварийного процесса, вектор РНМ, оказывается близок к границе зоны срабатывания, и ввиду использования классических ТТ имеет описанные ранее колебания [24]\commentnote{DISS-005-C0087}{Евдаков Алексей}{А.Е. Евдаков, А.А. Яблоков, Г.А. Филатова Исследование применения катушек Роговского для устройств БАВР. Материалы конференции «Радиоэлектроника , электротехника и энергетика», 29 февраля -- 2 марта 2024 г., Москва, с. 1194}, а при насыщении~-- более ярко выраженные (сигнал CT\_real). Стоит отметить, что далеко не у всех производителей устройств БАВР имеется гистерезис угловой характеристики, а также не все службы наладки задают нулевые выдержки времени на срабатывание и возврат органа РНМ. Ввиду этого данные колебания могут оказывать критическое влияние на поведение всей логики функционирования БАВР. Пример поведения различных векторов РНМ для такого аварийного процесса (Рисунок~\ref{fig:ch02_hodographs_external_fault}).

% DISS-005-B1320
Как мы можем видеть, использование катушек Роговского (сигнал RC) позволяет обеспечить наибольшую устойчивость функционирования в сравнении с классическими ТТ без насыщения и, особенно, при насыщении. При этом, стоит отметить, что ввиду физических особенностей кривая вектора РНМ для классических ТТ проходит область вблизи начала координат, что в некоторых случаях может даже привести к ускорению отпадания РНМ. Однако это требует более глубоких исследований функционирования устройства БАВР в целом, так как пусковые сигналы от других органов (27/27S и пуск по углу) могут ещё отсутствовать. Поэтому, если колебания вектора, обусловленные искажениями от апериодической составляющей не критичны (годографы для ТТ без насыщения), то может быть актуален расчёт вектора РНМ на основе восстановленного исходного тока путём интегрирования сигнала, полученного с КР.

% DISS-005-B1321
Годограф фазного РНМ на основе КР (PDR4) не приводится, так как кривая движения вектора аналогична с прямой последовательностью за исключением конечных точек.

% DISS-005-B1322
Применение усреднённого (CT\_RC) и гибридного (CT\_RC\_G) сигналов позволяет снизить влияние эффекта насыщения. При этом гибридный алгоритм значительно уменьшает данный колебательный процесс. Стоит отметить, что далеко не у всех производителей устройств БАВР имеется гистерезис угловой характеристики, а также не все службы наладки задают нулевые выдержки времени на срабатывание и возврат органа РНМ. Колебания вектора РНМ вблизи области срабатывания могут критически влиять на работу БАВР.

\begin{figure}[htbp]
  \centering
  \begin{tabular}{cc}
    % DISS-005-B1323
    \includegraphics[width=0.46\textwidth]{image69.png} &
    % DISS-005-B1324
    \includegraphics[width=0.46\textwidth]{image70.png} \\
    % DISS-005-B1325
    а) &
    % DISS-005-B1326
    б) \\
    % DISS-005-B1327
    \includegraphics[width=0.46\textwidth]{image71.png} &
    % DISS-005-B1328
    \includegraphics[width=0.46\textwidth]{image72.png} \\
    % DISS-005-B1329
    в) &
    % DISS-005-B1330
    г)
  \end{tabular}
  % DISS-005-B1331
  % DISS-005-B1332
  % DISS-005-B1333
  \caption{Годографы векторов РНМ при внешнем коротком замыкании:\\
  а) ПП для ТТ без насыщения, б) для ТТ с насыщением и КР,\\
  в) фазный для ТТ с насыщением, г) для ТТ с насыщением и КР с усреднением или гибридным алгоритмами (G),}
  \label{fig:ch02_hodographs_external_fault}
\end{figure}

% DISS-005-B1334
Второй рассматриваемый режим~-- это двойное замыкание на землю в разных точках сети 10~кВ, которое развилось из однофазного замыкания на землю (ОЗЗ). Сначала было замыкание фазы С на отходящей линии (точка К-2), которое перешло в двухфазное на землю за счёт замыкания фазы А вблизи секции шин на участке от силового трансформатора до выключателя и датчиков тока (К-3). Переходное сопротивление дуги моделируется дополнительным сопротивлением (в блоке «Fault» Simulink). Данный тип развития ОЗЗ в двухфазное замыкание является частым явлением. Пример поведения различных векторов РНМ для такого аварийного процесса (Рисунок~\ref{fig:ch02_hodographs_internal_fault}).

% DISS-005-B1335
В данном случае, аварию стоит расценивать как внутреннюю, и РНМ должен блокировать БАВР. Для случаев с реальными (насыщаемыми~-- сигнал CT\_real) ТТ, траектория вектора РНМ на начальном участке, где ещё не произошло насыщение, совпадает с идеальной (сигнал CT\_ideal). Однако после насыщения, она меняет своё направление и на короткий промежуток выходит из зоны по амплитуде (пересечение границы по окружности), что может привести к ложному срабатыванию БАВР. 

% DISS-005-B1336
Траектория векторов фазных реле направления мощности для каждой из фаз различна. Как видно, фаза А и фаза С (ввиду насыщения) выходят из зоны срабатывания. Так как логика их функционирования построена по условию «И» (то есть РНМ выдаёт сигнал о срабатывании, когда все вектора находятся в сработанном состоянии)~-- то это приведёт к «отпусканию» РНМ, с последующим ложным срабатывание БАВР. Возможность возникновения данных событий является причиной, по которой рекомендуется не использовать фазный РНМ в сетях 6-35~кВ. Поэтому указанный алгоритм реле направления мощности применяется лишь в сетях с несимметричной нагрузкой~-- например, 0,4~кВ.

% DISS-005-B1337
Применение усредненного и гибридного сигналов снижает колебания, вызванные эффектом насыщения ТТ, особенно гибридный алгоритм. Это позволяет избежать ложного срабатывания при внутреннем замыкании.

\begin{figure}[htbp]
  \centering
  \begin{tabular}{cc}
    % DISS-005-B1338
    \includegraphics[width=0.46\textwidth]{image73.png} &
    % DISS-005-B1339
    \includegraphics[width=0.46\textwidth]{image74.png} \\
    % DISS-005-B1340
    а) &
    % DISS-005-B1341
    б) \\
    % DISS-005-B1342
    \includegraphics[width=0.46\textwidth]{image75.png} &
    % DISS-005-B1343
    \includegraphics[width=0.46\textwidth]{image76.png} \\
    % DISS-005-B1344
    в) &
    % DISS-005-B1345
    г)
  \end{tabular}
  % DISS-005-B1346
  % DISS-005-B1347
  % DISS-005-B1348
  \caption{Годографы векторов РНМ при внутреннем коротком замыкании:\\
  а) ПП для ТТ без насыщения, б) для ТТ с насыщением и КР,\\
  в) фазный для ТТ с насыщением, г) для ТТ с насыщением и КР с усреднением или гибридным алгоритмами (G),}
  \label{fig:ch02_hodographs_internal_fault}
\end{figure}

% DISS-005-B1349
Результаты исследований в сложных режимах, в которых возможно некорректное срабатывание БАВР с ТТ, продемонстрировали актуальность применения КР для повышения стабильности и быстродействия. В связи с чем, было принято решение провести исследования на реальном энергообъекте с подключением БАВР с помощью КР. Одновременное же использование сигналов и с ТТ, и с КР считаем полезным и актуальным для будущих разработок.
